\documentclass[11pt,letterpaper]{article}

% -------------------------------------------------------------
% Packages
% -------------------------------------------------------------
\usepackage[margin=1in]{geometry}
\usepackage{amsmath,amssymb,amsthm}
\usepackage{mathrsfs}
\usepackage{graphicx}
\usepackage{hyperref}
\usepackage{parskip} % For better spacing between paragraphs
\usepackage{enumitem} % For customizing lists
\usepackage{listings}


% -------------------------------------------------------------
% Title and Author
% -------------------------------------------------------------
\title{\textbf{Cheat Sheet}\\
\large \textbf{An Introduction to the Analysis of Algorithms, Second Edition}\\
\small By Robert Sedgewick and Philippe Flajolet}
\author{\textit{Compiled by: Your Name}}
\date{\today}

% -------------------------------------------------------------
% New Commands
% -------------------------------------------------------------
\newcommand{\OGF}{\mathrm{OGF}}
\newcommand{\EGF}{\mathrm{EGF}}
\newcommand{\Prob}{\mathbb{P}}
\newcommand{\Expect}{\mathbb{E}}
\newcommand{\Var}{\mathrm{Var}}

\usepackage{listings}
\lstset{
  language=Java,
  basicstyle=\ttfamily,
  numbers=left,
  numberstyle=\small,
  stepnumber=1,
  numbersep=10pt,
  frame=single,
  captionpos=b,
  breaklines=true,
  breakatwhitespace=false,
  showstringspaces=false,
  tabsize=2,
}


\begin{document}

\maketitle
\tableofcontents
\clearpage

\section{Analysis of Algorithms}

\subsection{Why Analyze an Algorithm?}

\begin{itemize}
  \item \textbf{Efficient Resource Use:} Determine time and space requirements to see if an algorithm meets performance needs.
  \item \textbf{Basis for Comparison:} Compare different algorithms for the same task (e.g., sorting) under similar assumptions.
  \item \textbf{Scalability Insights:} Understand how performance grows with input size (\(N\)) using asymptotic notation.\footnote{Refer to the \emph{Notation} section of the document for definitions of Big-Oh, Big-Theta, etc.}
  \item \textbf{Practical Engineering:} Informs decisions on tradeoffs (e.g., time vs.\ memory, or time vs.\ battery usage on mobile devices).
  \item \textbf{Implementation Awareness:} A thorough analysis can highlight improvements in data layout or hardware exploitation.
\end{itemize}

\noindent
\emph{Real-world tip:} In modern software systems, analyzing an algorithm can prevent downstream issues with large data sets. For instance, if you anticipate handling multi-gigabyte logs or millions of records, an \(O(N^2)\) approach becomes impractical.

%-----------------------------------------------------
\subsection{Theory of Algorithms}

\begin{itemize}
  \item \textbf{Purpose:} Classify and bound performance.
  \item \textbf{Asymptotic Notation:} 
    \begin{itemize}
      \item \(O(\cdot)\) (upper bound), 
      \item \(\Omega(\cdot)\) (lower bound),
      \item \(\Theta(\cdot)\) (matching bound).
    \end{itemize}
  \item \textbf{Worst-Case Focus:} Provides a solid performance guarantee, though may be overly pessimistic.
  \item \textbf{Optimality:} Some algorithms match known lower bounds; for example, mergesort proves \(O(N \log N)\) is optimal for \emph{comparison}-based sorting.\footnote{See also Section 1.6 for “Why mergesort is optimal in worst-case scenarios.”}
\end{itemize}

\noindent
\emph{Cross-reference:} The concept of \emph{divide-and-conquer} recurrences is expanded in Chapter~2.  (See especially \S2.6.)

%-----------------------------------------------------
\subsection{Analysis of Algorithms}

\begin{itemize}
  \item \textbf{Exact vs.\ Approximate Models:} In practice, approximate (e.g., ignoring constant factors) can suffice for large-scale predictions.
  \item \textbf{Average-Case vs.\ Worst-Case:} Many practical applications benefit more from average-case results (e.g., expected running time on “typical” input).
  \item \textbf{Implementation Details:} Profiling often pinpoints the “hot spots,” making a more refined analysis possible.
  \item \textbf{Empirical Check:} Compare derived formulas against real measurements; if they match closely, confidence in the analysis grows.
\end{itemize}

\noindent
\emph{Lesson learned:} Iterative refinement of the analysis---backed by small empirical checks---often reveals hidden inefficiencies, such as poor data access patterns, that strictly mathematical models might gloss over.

%-----------------------------------------------------
\subsection{Average-Case Analysis}

\begin{itemize}
  \item \textbf{Input Models:} Often assume random permutations or random uniform distribution, especially for sorting or searching.
  \item \textbf{Significance:} Yields more accurate estimates in real-world scenarios where “adversarial” inputs (worst-case) rarely occur.
  \item \textbf{Tools:} Probability distributions, expected values, recurrences with average-case or probabilistic components.
  \item \textbf{Example:} Many sorts (including quicksort) have \(O(N \log N)\) average-case time, but simpler algorithms might degrade to \(O(N^2)\) in the worst case.
\end{itemize}

\noindent
\emph{Cross-reference:} For deeper analysis of specific average-case recurrences, see Chapter~2 (\S2.7) and the generating function techniques in Chapter~3.

%-----------------------------------------------------
\subsection{Example: Analysis of Quicksort}

\begin{itemize}
  \item \textbf{Partition-based Sorting:} Recursively partitions the array around a pivot.
  \item \textbf{Average Comparisons:} Roughly \(2N \ln N\).\footnote{In more detailed form, \(\approx 2(N+1)(H_{N+1}-3/2)\). See the document for exact derivation.}
  \item \textbf{Key Implementation Detail:} Choice of pivot influences performance. A “randomized” pivot often yields robust behavior.
  \item \textbf{Practical Tip:} Switching to an alternative sort (e.g., insertion sort) for small subarrays can improve constant factors.
  \item \textbf{Further Optimization:} Median-of-three pivot selection, sampling-based pivot selection, or so-called \emph{radix-exchange} variants if keys are known to be bit strings.
\end{itemize}

\noindent
\textbf{Code Example: Quicksort in Java}
\label{quicksort-prog}
\begin{lstlisting}[language=Java]
private void quicksort(int[] a, int lo, int hi) {
    if (hi <= lo) return;
    int i = lo - 1, j = hi;
    int v = a[hi];
    while (true) {
        while (a[++i] < v) ;        
        while (v < a[--j]) if (j == lo) break;
        if (i >= j) break;
        swap(a, i, j);
    }
    swap(a, i, hi);
    quicksort(a, lo, i - 1);
    quicksort(a, i + 1, hi);
}
\end{lstlisting}

\noindent
\emph{Lesson learned:} Randomizing pivot selection (e.g., picking a random index in \([\text{lo}..\,\text{hi}]\)) before partitioning ensures inputs effectively appear random, circumventing the worst-case pattern and making the average-case analysis predictions hold in practice.

%-----------------------------------------------------
\subsection{Asymptotic Approximations}

\begin{itemize}
  \item \textbf{Notation:} \(\sim\) denotes asymptotic equivalence; e.g.\ \(n!\sim \sqrt{2\pi n}\,(n/e)^n\).
  \item \textbf{Purpose:} Offer concise forms for large \(n\). Example: Harmonic number \(H_n \approx \ln n + \gamma\).
  \item \textbf{Efficiency:} Helps communicate “dominant term” and relevant constant factors for large inputs.
  \item \textbf{Cross-Reference:} Chapter~4 introduces the Euler--Maclaurin Summation Formula for analyzing sums and integrals in detail.
\end{itemize}

%-----------------------------------------------------
\subsection{Distributions}

\begin{itemize}
  \item \textbf{Beyond the Mean:} Standard deviation, variance, and higher moments often clarify the likelihood of \emph{typical} vs.\ extreme outcomes.
  \item \textbf{Limiting Distributions:} Many quantities in the analysis of algorithms approach well-known distributions (e.g., normal, Poisson).
  \item \textbf{Example:} Quicksort’s cost distribution becomes sharply peaked around its mean for large \(N\).
  \item \textbf{Empirical Check:} Histograms can validate that 90\%+ of runs fall near the predicted mean, confirming the “tightness” of the distribution.
\end{itemize}

\noindent
\emph{Practical note:} Systematic anomalies (e.g., partially sorted data, repeated keys) may shift distributions unexpectedly. Conduct limited randomized pre-processing (\texttt{shuffle}) to restore typical performance.

%-----------------------------------------------------
\subsection{Randomized Algorithms}

\begin{itemize}
  \item \textbf{Key Idea:} Inject randomness internally (e.g., random pivot selection in quicksort) rather than rely on random external inputs.
  \item \textbf{Predictable Performance:} Yields average-case behavior \emph{with high probability}, even if the input is adversarial.
  \item \textbf{Broader Applications:} Probabilistic data structures (Bloom filters), Monte Carlo methods, and more.
  \item \textbf{Cross-Reference:} For additional randomized methods, consult references and see advanced treatments in research on average-case complexity.
\end{itemize}

\noindent
\emph{Lesson learned:} Randomized algorithms reduce the risk that “bad” inputs degrade performance severely. The overhead of generating random numbers is often negligible compared to the gains.

%-----------------------------------------------------
% Code Example: Mergesort
%-----------------------------------------------------
\subsection{Example Code: Mergesort (Contrast with Quicksort)}

\begin{lstlisting}[language=Java, label=mergesort-prog]
private void mergesort(int[] a, int lo, int hi) {
    if (hi <= lo) return;
    int mid = lo + (hi - lo)/2;
    mergesort(a, lo, mid);
    mergesort(a, mid+1, hi);

    // Merge step
    int i = lo, j = mid+1;
    for (int k = lo; k <= hi; k++) aux[k] = a[k];
    for (int k = lo; k <= hi; k++) {
        if      (i > mid)          a[k] = aux[j++];
        else if (j > hi)           a[k] = aux[i++];
        else if (aux[j] < aux[i])  a[k] = aux[j++];
        else                       a[k] = aux[i++];
    }
}
\end{lstlisting}

\noindent
\emph{Comparison to Quicksort:}
\begin{itemize}
  \item \textbf{Worst-Case} mergesort is always \(O(N \log N)\).
  \item \textbf{Constant Factors} for mergesort are generally higher than for quicksort in practice.
  \item \textbf{Stable Sort:} Mergesort preserves item order for equal keys, which is sometimes crucial (e.g., stable merges of data).
  \item \textbf{Large Data Tip:} Mergesort excels in external sorting scenarios (e.g., on-disk merges).
\end{itemize}

%-----------------------------------------------------
% Summary for Chapter 1
%-----------------------------------------------------
\subsection{Summary}
\begin{itemize}
  \item Analyzing algorithms provides actionable insight for selecting, improving, or comparing methods.
  \item Worst-case performance bounds give theoretical guarantees; average-case analyses often align more closely with real-world performance.
  \item Randomized algorithms naturally yield predictable, near-average performance for a broad range of inputs.
  \item Mergesort and Quicksort are exemplars of two major sorting approaches; both highlight how recursive structures translate to mathematical recurrences.\footnote{See the document’s Recurrence Relations in Chapter~2 for further specifics on solution techniques.}
\end{itemize}

\noindent
\emph{Forward Reference:} 
\begin{itemize}
  \item \textbf{Chapter 2 (Recurrences)} delves deeper into how mergesort, quicksort, and other divide-and-conquer methods lead to solvable equations.
  \item \textbf{Chapters 3--4} introduce generating functions and asymptotic expansions for analyzing sums and series more systematically.
\end{itemize}

\clearpage

%=====================================================
% Chapter 2: Recurrence Relations
%=====================================================
\section{Recurrence Relations}

\subsection{Basic Properties}
\begin{itemize}
  \item \textbf{Definition:} A recurrence expresses \(a_n\) in terms of values at smaller indices (e.g., \(a_{n-1}, a_{n-2}\)).
  \item \textbf{Motivation:} Many algorithms (e.g., \emph{divide-and-conquer}) produce a direct recursive decomposition of cost or structure that leads to a recurrence.
  \item \textbf{Key Insight:} Recurrences capture both \emph{exact} and \emph{approximate} descriptions of algorithmic costs.
  \item \textbf{Terminology/Classification:}
  \begin{itemize}
    \item \emph{Order}: Number of previous terms (first-order, second-order, etc.).
    \item \emph{Linearity}: If the right-hand side is a linear combination of earlier terms, it is \emph{linear}.
    \item \emph{Constant vs.\ variable coefficients}.
    \item \emph{Homogeneous vs.\ inhomogeneous}.
  \end{itemize}
  \item \emph{Actionable Tip:} Always compute small values to test growth and confirm patterns before seeking a closed-form solution.
  \item \emph{Cross-Reference:} For an example of recurrences from sorting, see \hyperref[quicksort-proof]{Quicksort Analysis} in Chapter~1.
\end{itemize}

%-----------------------------------------------------
\subsection{First-Order Recurrences}
\begin{itemize}
  \item \textbf{Basic Forms:}
    \begin{itemize}
      \item \emph{Product-like}: \(a_n = x_n \, a_{n-1} \implies a_n = \prod_{k=1}^n x_k\).
      \item \emph{Sum-like}: \(a_n = a_{n-1} + y_n \implies a_n = \sum_{k=1}^n y_k\).
    \end{itemize}
  \item \textbf{Iteration Method}: Directly “unroll” the recurrence to reduce to known sums/products.
  \item \textbf{Example:} 
    \[
      a_n = 2\,a_{n-1} + 1
      \quad\Longrightarrow\quad
      \text{Divide both sides by } 2^n \text{ and iterate.}
    \]
  \item \textbf{Summation Factor Technique:} Multiply each side by a term (e.g.\ \(\frac{1}{2^n}\)) so that the left-hand side telescopes when iterating.
  \item \emph{Cross-Reference:} See \S2.5 for advanced solution methods, including \emph{change of variable}, \emph{repertoire}, and \emph{perturbation}.
\end{itemize}

%-----------------------------------------------------
\subsection{Nonlinear First-Order Recurrences}
\begin{itemize}
  \item \textbf{Example Types:} 
    \[
      a_n = 1/(1 + a_{n-1}), 
      \quad a_n = \sin(a_{n-1}), 
      \quad a_n = \alpha\, a_{n-1}^2 + \ldots
    \]
  \item \textbf{Common Behaviors}:
    \begin{itemize}
      \item \textit{Simple Convergence}: \(a_n \to \alpha\) under iteration if \(|f'(\alpha)| < 1\).
      \item \textit{Quadratic Convergence}: E.g.\ Newton’s method can \emph{double} the number of correct digits each iteration.
      \item \textit{Slow Convergence}: If \(|f'(\alpha)| = 1\), the sequence \(a_n\) may converge slowly.
    \end{itemize}
  \item \emph{Actionable Tip:} Check \(\lim_{n \to \infty} a_n\) by solving \(\alpha = f(\alpha)\). Numerically verify early terms to confirm behavior.
\end{itemize}

%-----------------------------------------------------
\subsection{Higher-Order Recurrences}
\begin{itemize}
  \item \textbf{Linear Recurrences with Constant Coefficients}: 
    \[
      a_n = x_1\,a_{n-1} + x_2\,a_{n-2} + \cdots + x_t\,a_{n-t}.
    \]
    \begin{itemize}
      \item \emph{Characteristic Polynomial} \(q(z) = z^t - x_1 z^{t-1} - \cdots - x_t.\)
      \item Solutions combine terms of the form \(n^j \beta^n\), where \(\beta\) is a root (possibly repeated) of \(q\).
    \end{itemize}
  \item \textbf{Fibonacci Example}:
    \[
      F_n = F_{n-1} + F_{n-2},\quad F_0=0,\;F_1=1
      \;\Longrightarrow\;
      F_n \approx \frac{\phi^n}{\sqrt{5}}\quad(\phi\approx 1.618).
    \]
  \item \emph{Tip for Implementation:} Use the closed form \emph{only} if you need exact integer values for smaller \(n\). Otherwise, asymptotic approximations or generating functions (see Chapter~3) are often easier to manipulate.
\end{itemize}

%-----------------------------------------------------
\subsection{Methods for Solving Recurrences}
\begin{itemize}
  \item \textbf{Change of Variable}: Simplify a complex form by substituting \(b_n = g(a_n)\). If you can map it to a simpler linear recurrence, you can solve it exactly.
  \item \textbf{Repertoire (Method of Undetermined Coefficients)}: Guess candidate solutions (like \(n\), \(n^2\), \(\log n\)) to build a “library” of partial solutions, then linearly combine them to solve the target recurrence.
  \item \textbf{Bootstrapping}:
    \begin{enumerate}
      \item Guess approximate growth (e.g., \(a_n \sim c \cdot n \log n\)).
      \item Substitute back into the recurrence.
      \item Refine constants or extra terms until stable.
    \end{enumerate}
  \item \textbf{Perturbation}: Compare your recurrence to a simpler one for which you have a known solution (e.g., drop lower-order terms) and bound the difference carefully.
\end{itemize}

%-----------------------------------------------------
\subsection{Binary Divide-and-Conquer Recurrences and Binary Numbers}
\begin{itemize}
  \item \textbf{Typical Form}:
    \[
      C_N = C_{\lfloor N/2 \rfloor} + C_{\lceil N/2 \rceil} + f(N)
      \quad\Longrightarrow\quad C_N \text{ often } \sim N \log N.
    \]
  \item \textbf{Examples}:
    \begin{itemize}
      \item \emph{Binary Search}: \(B_N = B_{\lfloor N/2 \rfloor} + 1\), solution \(B_N = \lfloor \log_2 N \rfloor + 1.\)
      \item \emph{Mergesort}: \(C_N = C_{\lfloor N/2\rfloor} + C_{\lceil N/2\rceil} + N.\) Exact solution:
        \[
          C_N = N \lfloor \log_2 N \rfloor + 2N - 2^{\lfloor \log_2 N\rfloor +1}.
        \]
      \item \emph{Other} “Equal-split” or “nearly-equal-split” methods similarly yield \(\Theta(N \log N)\).
    \end{itemize}
  \item \textbf{Connection to Binary Representation}: 
    \begin{itemize}
      \item \(\lfloor \log_2 N \rfloor\) measures the number of bits needed.
      \item Splitting “in half” often maps to shifting the binary representation by 1 bit.
      \item Recurrences can exhibit periodic or fractal-like behavior (see figures of \(C_N - N \log N\)).
    \end{itemize}
  \item \emph{Practical Tip:} In analysis, you often only need an asymptotic \(\Theta(\cdot)\) argument, but for fine-tuned performance comparisons, be mindful of integer \(\lfloor \cdot \rfloor\) and \(\lceil \cdot \rceil\) effects.
\end{itemize}

\subsection{Code Example: Computing Recurrence Values in Practice}
\label{prog:recurrence}
\begin{lstlisting}[language=Java]
double[] C = new double[maxN + 1];
C[0] = 0.0;
for (int N = 1; N <= maxN; N++) {
    C[N] = N + 1.0; // add f(N) = N + 1
    for (int k = 0; k < N; k++) {
        C[N] += (C[k] + C[N - 1 - k]) / N;
    }
}
// C[N] now holds the solution for each 1 <= N <= maxN
\end{lstlisting}

\noindent
\emph{Actionable Insight:} A direct iterative or memoized approach is often preferable to a naive recursive approach, which can explode in complexity by recomputing the same values multiple times.

%-----------------------------------------------------
\subsection{General Divide-and-Conquer Recurrences}
\begin{itemize}
  \item \textbf{Master-Style Form:} 
    \[
      a(N) = \alpha\,a\!\Bigl(\frac{N}{\beta}\Bigr) + f(N),
      \quad \beta > 1.
    \]
  \item \textbf{Three Cases for \(\alpha\) vs.\(\beta\)} (assuming \(f(N) \approx N^\gamma(\log N)^\delta\)):
    \begin{enumerate}
      \item \(\gamma < \log_\beta \alpha\): \(\quad a(N) = \Theta\bigl(N^{\log_\beta \alpha}\bigr)\)
      \item \(\gamma = \log_\beta \alpha\): \(\quad a(N) = \Theta\bigl(N^\gamma (\log N)^{\delta+1}\bigr)\)
      \item \(\gamma > \log_\beta \alpha\): \(\quad a(N) = \Theta\bigl(N^\gamma (\log N)^\delta\bigr)\)
    \end{enumerate}
  \item \textbf{Example:} Mergesort has \(\alpha=2\), \(\beta=2\), and \(f(N)=N\). Then \(\gamma=1\), \(\log_\beta \alpha = 1\). So mergesort’s cost is \(\Theta(N \log N)\).
  \item \emph{Guidance:} This approach is often used in theoretical algorithm analysis to get big-picture complexity bounds. For more precise results, combine with exact recurrences or generating functions.
\end{itemize}

%-----------------------------------------------------
% Summary for Chapter 2
%-----------------------------------------------------
\subsection{Summary}
\begin{itemize}
  \item Recurrences arise naturally from recursive/iterative program structures.
  \item Iteration, summation factors, and change of variables are critical tools for \textbf{first-order} recurrences.
  \item \textbf{Higher-order linear recurrences} with constant coefficients have exact solutions via the characteristic polynomial (roots \(\rightarrow\) terms \(\beta^n\)).
  \item \textbf{Nonlinear or variable-coefficient recurrences} can often be approximated or transformed, but exact closed forms may not exist.
  \item \textbf{Divide-and-conquer recurrences} capture major classes of algorithms and exhibit \(\Theta(n\log n)\) behavior in many balanced scenarios (e.g.\ mergesort).
  \item \emph{Real-World Tip:} Validate your formula for small \(n\) using a direct or memoized computation, then use closed-form or asymptotic approximations for large \(n\).
  \item \textbf{Forward Reference:} Chapter~3 shows how \emph{generating functions} solve many recurrences systematically, often bypassing messy ad hoc manipulations.
\end{itemize}
\clearpage


%=====================================================
% Chapter 3: Generating Functions
%=====================================================
\section{Generating Functions}

\subsection{Ordinary Generating Functions (OGFs)}
\begin{itemize}
  \item \textbf{Definition:} For a sequence \(\{a_k\}\), its \emph{ordinary generating function} (OGF) is
    \[
      A(z) \;=\;\sum_{k=0}^{\infty} a_k\,z^k.
    \]
  \item \textbf{Key Notation:} If \(A(z) = \sum_{k\ge0} a_k\,z^k\), then \([z^k]A(z) = a_k\).
  \item \textbf{Elementary OGFs:} 
    \[
      \frac{1}{1-z} = 1+z+z^2+\cdots,\quad
      \frac{1}{(1-z)^2} = 1+2z+3z^2+\cdots,\quad
      \frac{1}{1-2z} = \sum_{k\ge0} (2^k)\,z^k.
    \]
  \item \textbf{Standard Transformations:}
    \begin{itemize}
      \item \textit{Differentiation:} \(\frac{d}{dz}\bigl(\sum a_k z^k\bigr) = \sum k\,a_k\,z^{k-1}\).
      \item \textit{Multiplication by \((1-z)\)}: Convolution/differences of coefficients (cf.\ \S\ref{sec:conv-diff}).
    \end{itemize}
  \item \emph{Actionable Tip:} Use OGFs to solve linear recurrences by transforming them into rational functions, then recover coefficients via expansions or partial fractions.
  \item \emph{Cross-Reference:} See also \hyperref[sec:bivariate]{Bivariate Generating Functions} to track a second parameter (e.g.\ “cost”).
\end{itemize}

%-----------------------------------------------------
\subsection{Exponential Generating Functions (EGFs)}
\begin{itemize}
  \item \textbf{Definition:} For a sequence \(\{a_k\}\), its \emph{exponential generating function} (EGF) is
    \[
      A^*(z) \;=\;\sum_{k=0}^{\infty} a_k \,\frac{z^k}{k!}.
    \]
  \item \textbf{Key Notation:} If \(A^*(z) = \sum_{k\ge0} a_k\,\frac{z^k}{k!}\), then \(k!\,[z^k]A^*(z)=a_k\).
  \item \textbf{Usage:} Often arises when counting \emph{labelled} structures: dividing by \(k!\) accounts for permutations of labels.
  \item \textbf{Elementary EGFs:} 
    \[
      e^z = \sum_{k\ge0} \frac{z^k}{k!},\quad
      \frac{1}{1-z} = \sum_{k\ge0} z^k,\quad
      \sin(z),\;\cos(z),\;\ldots
    \]
  \item \textbf{Transformations:}
    \begin{itemize}
      \item \textit{Binomial Convolution:} If \(A^*(z)\) and \(B^*(z)\) are EGFs, their product \(\sum a_k \frac{z^k}{k!}\)\(\times\)\(\sum b_k \frac{z^k}{k!}\) yields a sequence with binomial-sum coefficients.
      \item \textit{Differentiation:} \(\frac{d}{dz}\Bigl(\sum a_k \frac{z^k}{k!}\Bigr) = \sum k\,a_k \frac{z^{k-1}}{(k-1)!}\).
    \end{itemize}
  \item \emph{Tip:} EGFs help in analyzing combinatorial classes where elements are distinct, e.g.\ permutations (\S\ref{sec:permutations}).
\end{itemize}

%-----------------------------------------------------
\subsection{Generating Functions for Solving Recurrences}
\begin{itemize}
  \item \textbf{Method:}
    \begin{enumerate}
      \item Multiply both sides of the recurrence by \(z^n\) (OGF) or \(\frac{z^n}{n!}\) (EGF).
      \item Sum over \(n\).
      \item Identify known forms or transform into partial fractions / known expansions.
      \item Invert to find coefficients.
    \end{enumerate}
  \item \textbf{Example (Simple):} 
    \[
      a_n \;=\; a_{n-1}+1,\;\; a_0=0
      \quad\Longrightarrow\quad
      A(z)=\frac{z}{(1-z)^2},\;\; a_n = n.
    \]
  \item \textbf{Example (Fibonacci):}
    \[
      F_n = F_{n-1}+F_{n-2},\;\;F_0=0,\;F_1=1 
      \quad\Longrightarrow\quad
      F(z)=\frac{z}{1-z-z^2}.
    \]
  \item \emph{Actionable Tip:} Automate partial-fraction expansions with symbolic tools for complex/large-order recurrences.
\end{itemize}

%-----------------------------------------------------
\subsection{Expanding Generating Functions}
\begin{itemize}
  \item \textbf{Taylor Series:} 
    \[
      f(z) = \sum_{n\ge0} \frac{f^{(n)}(0)}{n!}\,z^n
    \]
    is the formal power series expansion. 
  \item \textbf{Partial Fractions / Algebraic Decomposition:} 
    \[
      \frac{1}{1 - az - bz^2} \quad\mapsto\quad \text{roots + expansions.}
    \]
  \item \textbf{Binomial Theorem (Generalized)}:
    \[
      (1+z)^\alpha = \sum_{k\ge0} \binom{\alpha}{k} z^k,\quad
      \binom{\alpha}{k} = \frac{\alpha(\alpha-1)\cdots(\alpha-k+1)}{k!}.
    \]
  \item \textit{Lesson:} Many known series expansions (e.g.\ for \(\exp(z)\), \(\log(1-z)\), \((1 - z)^\alpha\)) suffice for most algorithmic recurrences.
\end{itemize}

%-----------------------------------------------------
\subsection{Transformations with Generating Functions}
\label{sec:conv-diff}
\begin{itemize}
  \item \textbf{Convolution:} 
    \[
      A(z)B(z)=\Bigl(\sum a_k z^k\Bigr)\Bigl(\sum b_k z^k\Bigr)
      =\sum_{n\ge0}\Bigl(\sum_{k=0}^n a_k b_{n-k}\Bigr) z^n.
    \]
  \item \textbf{Shift/Partial Sum:}
    \[
      \frac{A(z)}{1-z} = \sum_{n\ge0} \Bigl(\sum_{k=0}^n a_k\Bigr)z^n.
    \]
  \item \textbf{Differentiation:} 
    \[
      \frac{d}{dz}A(z) = \sum_{n\ge1} n\,a_n\,z^{n-1}.
    \]
  \item \emph{Real-World Tip:} Convolution identities often yield combinatorial identities (e.g.\ Vandermonde’s Convolution) directly.
\end{itemize}

%-----------------------------------------------------
\subsection{Functional Equations on Generating Functions}
\begin{itemize}
  \item \textbf{Linear Examples}: 
    \[
      f(z)=c + \alpha z f(z) + \beta z^2 f(z),\;\;
      \Longrightarrow\; f(z) = \dots
    \]
  \item \textbf{Nonlinear Examples}: 
    \[
      T(z)=1+z\,T(z)^2 \quad(\text{Catalan trees}),\quad 
      T(z)=z\,e^{T(z)} \quad(\text{labelled trees}).
    \]
  \item \textbf{Differential Examples}:
    \[
      z\,\frac{d}{dz}Q(z)=2\,\frac{z}{(1-z)^3}+ \dots 
      \quad\Longrightarrow\quad \text{(Quicksort analysis)}
    \]
  \item \emph{Iterative Approach:} Repeated substitution often reveals structure or singularity. 
  \item \emph{Cross-Reference:} See \S\ref{sec:quicksort-med3} for the “median-of-three quicksort” functional equation approach.
\end{itemize}

%-----------------------------------------------------
\subsection{Solving the Quicksort Median-of-Three Recurrence}
\label{sec:quicksort-med3}
\begin{itemize}
  \item \textbf{Recurrence} (from Chapter 1):
    \[
      C_N = N+1 + \sum_{k=0}^{N-1}\binom{N}{k}\frac{1}{2^N}\bigl(C_k + C_{N-1-k}\bigr).
    \]
  \item \textbf{Generating Function:} 
    \[
      C(z)= \sum_{N\ge0} C_N\,z^N 
      \;\Longrightarrow\;\text{Derive a high-order differential eq.}
    \]
  \item \textbf{Technique:} Multiply, sum, factor polynomial operators (Euler-type differential equation), partial fractions or expansions.
  \item \textbf{Result (High Level):} 
    \[
      C_N \approx 2N\ln N \quad(\text{constant factor differs slightly from standard quicksort}).
    \]
  \item \emph{Actionable Tip:} Larger pivot samples yield diminishing returns; the GF method quantifies these trade-offs precisely.
\end{itemize}

%-----------------------------------------------------
\subsection{Counting with Generating Functions}
\begin{itemize}
  \item \textbf{Method:} “Each object contributes one to the coefficient of \(z^{\text{(size)}}\).” 
  \item \textbf{Example (Binary Trees)}: 
    \[
      T(z)=1+\sum_{t_L,t_R\in T}z^{1+|t_L|+|t_R|}
      =1+z\,T(z)^2,
      \quad T(z)=\frac{1-\sqrt{1-4z}}{2z}.
    \]
    Coefficients are \emph{Catalan} numbers.
  \item \textbf{Coin Change Problem}:
    \[
      D(z)=\frac{1}{(1-z)(1-z^5)(1-z^{10})(1-z^{25})\cdots}
      \quad\Longrightarrow\quad[z^n]D(z)\;=\;\text{count of ways}.
    \]
  \item \emph{Lesson:} GFs offer a systematic approach to combinatorial enumeration: each structural choice \(\leftrightarrow\) a product or convolution in the GF.
\end{itemize}

%-----------------------------------------------------
\subsection{Probability Generating Functions (PGFs)}
\begin{itemize}
  \item \textbf{Definition:} For a discrete r.v.\ \(X\), \(\text{PGF}(u)=\mathbb{E}[u^X] = \sum_{k\ge0}\Pr(X=k)\,u^k.\)
  \item \textbf{Key Use:} 
    \[
      \mathbb{E}[X]=\text{PGF}'(1), 
      \quad \mathrm{Var}(X)=\text{PGF}''(1)+\text{PGF}'(1)-[\text{PGF}'(1)]^2.
    \]
  \item \textbf{Example (Binomial):}
    \[
      P(u) = (q + p\,u)^N,\quad \mathbb{E}[X]=pN,\quad \mathrm{Var}(X)=pq\,N.
    \]
  \item \textbf{Contrast with BGFs:} PGFs track cost distribution directly; BGFs in \S\ref{sec:bivariate} often prove more flexible for analyzing structure \emph{and} cost together.
\end{itemize}

%-----------------------------------------------------
\subsection{Bivariate Generating Functions (BGFs)}
\label{sec:bivariate}
\begin{itemize}
  \item \textbf{Definition:} For a doubly indexed sequence \(\{a_{n,k}\}\):
    \[
      A(z,u)=\sum_{n,k\ge0} a_{n,k}\,z^n\,u^k.
    \]
  \item \textbf{Motivation:} Track two parameters (e.g.\ size \(n\), cost \(k\)). 
  \item \textbf{Cumulative Generating Function (CGF)}:
    \[
      \frac{\partial}{\partial u}A(z,u)\Bigl|_{u=1} \quad\text{captures total/average cost.}
    \]
  \item \textbf{Mean Calculation:} If \(p_{n}(u)=\sum_k a_{n,k}u^k\) enumerates all size-\(n\) objects, then
    \[
      \text{cumulated cost} = p_n'(1),\quad
      \text{average cost} = \frac{p_n'(1)}{p_n(1)}.
    \]
  \item \textbf{Example (Binary Strings):}
    \[
      B(z,u)=\sum_{n\ge0}\sum_{k\ge0}\binom{n}{k}z^n\,u^k=(1+u)^n z^n \dots
    \]
  \item \emph{Lesson:} BGFs unify combinatorial counting (via \(z^n\)) and cost distribution (via \(u^k\)) in a single function.
\end{itemize}

%-----------------------------------------------------
\subsection{Special Functions}
\begin{itemize}
  \item \textbf{Bernoulli Numbers/Polynomials}, \textbf{Stirling Numbers}, \textbf{Eulerian Numbers}, etc.
  \item \textbf{Why “Special”?} These sequences arise in so many contexts (combinatorial formulas, expansions, algorithmic recurrences, etc.) that they have been studied extensively.
  \item \textbf{Examples:}
    \[
      \text{Bernoulli: } \frac{z}{e^z-1}=\sum_{n\ge0}B_n\frac{z^n}{n!},\quad
      \text{Stirling (1st/2nd kind): expansions involving falling factorials.}
    \]
  \item \emph{Cross-Reference:} Used frequently in \hyperref[sec:asymptotic-ch4]{Chapter 4} (Asymptotic Approximations) for Euler--Maclaurin summation.
\end{itemize}

%-----------------------------------------------------
% Summary for Chapter 3
%-----------------------------------------------------
\subsection{Summary}
\begin{itemize}
  \item Generating functions (ordinary/exponential) systematically encode entire sequences in a single functional form.
  \item Solving recurrences via GF transformations is often more streamlined than ad hoc “unrolling.”
  \item Functional equations on GFs capture the recursive structure of many algorithms/data structures (\emph{binary trees}, \emph{quicksort}, etc.).
  \item Probability generating functions and bivariate generating functions highlight how to derive moments (mean, variance) \emph{and} count combinatorial objects.
  \item Special functions (e.g., Bernoulli numbers, Stirling numbers) appear naturally and are well-studied resources for expansions and approximations.
  \item \emph{Forward Reference:} In \hyperref[sec:asymptotic-ch4]{Chapter 4}, we use asymptotic expansions to extract leading terms from GFs more efficiently; in \hyperref[ch:analytic-comb]{Chapter 5}, we embrace them as a primary tool in \emph{Analytic Combinatorics}.
\end{itemize}
\clearpage

%=====================================================
% Chapter 4: Asymptotic Approximations
%=====================================================
\section{Asymptotic Approximations}

\subsection{Notation for Asymptotic Approximations}
\begin{itemize}
  \item \textbf{Big-O, little-o, and Tilde}:
    \begin{itemize}
      \item $g(N) = O(f(N)) \;\Longleftrightarrow\; \bigl|\tfrac{g(N)}{f(N)}\bigr|\text{ bounded}$.
      \item $g(N) = o(f(N)) \;\Longleftrightarrow\; \lim_{N\to\infty}\tfrac{g(N)}{f(N)} = 0$.
      \item $g(N) \sim f(N) \;\Longleftrightarrow\; \lim_{N\to\infty}\tfrac{g(N)}{f(N)} = 1$.
    \end{itemize}
  \item \textbf{Usage:} 
    \[
      g(N)=f(N) + O(h(N))\;\;\Rightarrow\;\;\text{error term at most const.}\times h(N).
    \]
  \item \textbf{Exponential and Polynomial Growth}: 
    \begin{itemize}
      \item If $\alpha < \beta$, then $\alpha^N$ is exponentially small relative to $\beta^N$.
      \item If a term is “exponentially small,” it is smaller than any polynomial term $n^{-m}$, for large $n$.
    \end{itemize}
  \item \emph{Actionable Tip:} Reserve $O(\cdot)$ for bounding true “error terms.” Writing $f(N) = O(g(N))$ with no constant factor can lose useful precision.
\end{itemize}

%-----------------------------------------------------
\subsection{Asymptotic Expansions}
\begin{itemize}
  \item \textbf{Concept}: An \emph{asymptotic series} for $f(N)$ is a sum $c_0g_0(N)+c_1g_1(N)+\dots$ such that truncating at any stage yields $f(N)= c_0g_0(N)+\cdots+c_k g_k(N) + o\bigl(g_k(N)\bigr).$
  \item \textbf{Examples}:
    \begin{itemize}
      \item Harmonic number $H_N = \ln N + \gamma + \frac{1}{2N} - \frac{1}{12N^2} + \cdots$.
      \item Factorial $\displaystyle N! \sim \sqrt{2\pi N}\,\bigl(\tfrac{N}{e}\bigr)^N\Bigl(1+\tfrac{1}{12N}+\cdots\Bigr)$ (Stirling's formula).
    \end{itemize}
  \item \textbf{Finite truncations} such as $f(N)=c_0g_0(N)+c_1g_1(N)+O\bigl(g_2(N)\bigr)$ are typical in algorithm analysis, e.g.\ quicksort comparisons $= 2N\ln N + (2\gamma-2)N + O(\ln N)$.
  \item \emph{Lesson:} Each extra term typically yields a more accurate approximation for large $N$. 
\end{itemize}

%-----------------------------------------------------
\subsection{Manipulating Asymptotic Expansions}
\begin{itemize}
  \item \textbf{Algebraic Rules with Big-O Notation}:
    \[
      O(A)+O(B)\to O(\max\{A,B\}),\quad O(A)\times O(B)\to O(AB),
      \quad \dots
    \]
  \item \textbf{Expansion Methods}:
    \begin{enumerate}
      \item \emph{Substitution}: apply known expansions (e.g.\ Taylor) with $x=\frac{1}{N},\,\ln(N-2)\sim\ln N$, etc.
      \item \emph{Factoring}: factor out large terms to isolate “error” portion.
      \item \emph{Multiplying/Dividing}: carefully combine expansions, preserving only terms up to the desired order.
      \item \emph{Exponentiation/Logarithm}: $f(N)=\exp\{\ln(f(N))\}$ transforms multiplication into addition, etc.
      \item \emph{Reversion}: invert series expansions $y=x+c_2x^2+\cdots$ to find $x$ as a series in $y$.
    \end{enumerate}
  \item \emph{Actionable Tip:} \emph{Always track the largest error term.} Smaller-order terms can be dropped, but be mindful of possible cancellations requiring extra precision.
\end{itemize}

%-----------------------------------------------------
\subsection{Asymptotic Approximations of Finite Sums}
\begin{itemize}
  \item \textbf{Sum vs.\ Integral}: $\sum_{k=a}^{b}f(k)$ can often be approximated by $\int_{a}^{b}f(x)\,dx$.
  \item \textbf{Tail Bounds}:
    \begin{itemize}
      \item Rapidly \emph{decreasing} summands $\implies$ infinite tail sum is negligible.
      \item Rapidly \emph{increasing} summands $\implies$ last term often dominates.
    \end{itemize}
  \item \textbf{Refined Approach}: \emph{Euler--Maclaurin} formula.
\end{itemize}

%-----------------------------------------------------
\subsection{Euler--Maclaurin Summation}
\begin{itemize}
  \item \textbf{Classical Formula} (simplified):
    \[
      \sum_{k=a}^{b-1} f(k)
      \;=\;\int_{a}^{b} f(x)\,dx \;+\;\frac{f(a)+f(b)}{2}
      \;+\;\sum_{i=1}^{m} \frac{B_{2i}}{(2i)!}\,\bigl(f^{(2i-1)}(b)-f^{(2i-1)}(a)\bigr)+R_m,
    \]
    where $B_{2i}$ are Bernoulli numbers and $R_m$ is a remainder.
  \item \textbf{Use Case 1 (Fixed Interval)}: Summation approximates integral as step size $\to 0$.
  \item \textbf{Use Case 2 (1 to $N$) “Discrete form”}:
    \[
      \sum_{k=1}^{N} f(k)=\int_{1}^{N}f(x)\,dx + Cf + \dots
      \quad(\text{where $Cf$ is Euler--Maclaurin constant}).
    \]
  \item \textbf{Applications}:
    \begin{itemize}
      \item Harmonic number $H_N = \ln N + \gamma + \frac{1}{2N}-\cdots$.
      \item $\ln N! = N\ln N - N + \frac12\ln(2\pi N)+\cdots$ (Stirling).
    \end{itemize}
  \item \emph{Lesson:} Let derivatives of $f(x)$ vanish or be bounded; each additional term in the expansion yields an extra $\frac{B_{2i}}{(2i)!}f^{(2i-1)}(N)$-type correction.
\end{itemize}

%-----------------------------------------------------
\subsection{Bivariate Asymptotics}
\begin{itemize}
  \item \textbf{Context}: Summands depend on two parameters (index $k$ and size $N$).
  \item \textbf{Strategy}: Identify range of $k$ where the summand is significant and apply expansions differently in “main region” vs.\ “tails” (Laplace method).
  \item \textbf{Examples}:
    \begin{itemize}
      \item \emph{Ramanujan Q-Function}: $\displaystyle Q(N)=\sum_{k=0}^{N}\prod_{j=0}^{k-1}\bigl(1-\tfrac{j}{N}\bigr).$
      \item \emph{Binomial Distribution}: $\displaystyle \binom{N}{k}\frac{1}{2^N}$.
      \item \emph{Poisson Distribution limit}: For $p=\lambda/N$, $\binom{N}{k}p^k(1-p)^{N-k}\to \frac{\lambda^k e^{-\lambda}}{k!}$.
    \end{itemize}
\end{itemize}

%-----------------------------------------------------
\subsection{Laplace Method}
\begin{itemize}
  \item \textbf{Concept}: Evaluate $\sum a_N(k)$ or $\int e^{M f(x)}dx$ by restricting to regions where the summand (or integrand) is largest.
  \item \textbf{Steps} (for discrete sums):
    \begin{enumerate}
      \item Truncate range to isolate significant terms.
      \item Use expansions in the main region of interest.
      \item Bound tails outside that region (often exponentially small).
    \end{enumerate}
  \item \textbf{Applications}:
    \begin{itemize}
      \item Ramanujan Q-Function $\sim \sqrt{\tfrac{\pi}{2}}\,\sqrt{N}$.
      \item Summation of binomial-like terms (e.g.\ normal approximation).
      \item “Catalan-type sums” leading to integrals with $x e^{-x^2/2}$, etc.
    \end{itemize}
\end{itemize}

%-----------------------------------------------------
\subsection{"Normal" Examples from the Analysis of Algorithms}
\begin{itemize}
  \item \textbf{Summations Resembling $\binom{N}{k}2^{-N}F(k)$:} 
    \[
      \sum_{k} \binom{N}{k}\,F(k)\,\tfrac{1}{2^N}.
    \]
  \item \textbf{Methods}: Normal approximation + Euler--Maclaurin, or direct Laplace argument on partial range + bounding tails.
  \item \textbf{Examples}:
    \begin{itemize}
      \item \emph{$2$-ordered sequences}, \emph{inversions}, \emph{Batcher’s odd-even mergesort} times.
      \item Typically: $\sum_k \binom{N}{k}\frac{1}{2^N} F(k)\approx \int F(x)\,\phi(\dots)\,dx$ for some Gaussian form $\phi$.
    \end{itemize}
\end{itemize}

%-----------------------------------------------------
\subsection{"Poisson" Examples from the Analysis of Algorithms}
\begin{itemize}
  \item \textbf{Contrasting with Normal Region:} 
    \[
      \text{Poisson-like sums: } \sum_{k\ge0} e^{-\lambda}\,\frac{\lambda^k}{k!}\,f(k).
    \]
  \item \textbf{Tries and Radix-Exchange Sort}:
    \begin{itemize}
      \item Summation with $\displaystyle \binom{N}{j}$ but small $j$ (exponentially decreasing for large $j$).
      \item Yields a main portion for $j \approx O(\ln N)$ and a tail portion negligible.
    \end{itemize}
  \item \textbf{Trie Sum Example}:
    \[
      S_N = \sum_{j\ge0}\Bigl(1-e^{-N/2^j}\Bigr)\sim \lg N.
    \]
\end{itemize}

%-----------------------------------------------------
% Summary for Chapter 4
%-----------------------------------------------------
\subsection{Summary}
\begin{itemize}
  \item \textbf{Asymptotics} is critical for “suppressing” messy details and comparing leading terms in algorithm analysis.
  \item \textbf{Basic Tools}:
    \begin{itemize}
      \item Big-O, small-o, and $\sim$ notations.
      \item Asymptotic expansions (e.g.\ expansions for $\ln N!$, $H_N$, $\binom{2N}{N}$).
      \item Euler--Maclaurin summation for sum/integral approximations.
      \item Laplace method for restricting to “peak” regions or bounding tails.
    \end{itemize}
  \item \textbf{Key “Normal”/“Poisson” Approximations}: 
    \begin{itemize}
      \item Binomial distribution $\sim$ normal near center (bell curve), $\sim$ Poisson if $p=\lambda/N$.
      \item Ramanujan $Q,R$-distributions $\sim$ partial sums of exponentials, leading to $\sqrt{\pi N/2}$ forms.
    \end{itemize}
  \item \textbf{Insights for Algorithms}:
    \begin{itemize}
      \item For sums with factorial or binomial coefficients, expansions give precise average-case results.
      \item Logarithmic/exponential phenomena frequently appear (Mergesort $\sim N\log N$, tries $\sim N\log N$, etc.).
    \end{itemize}
  \item \emph{Forward Reference:} In \hyperref[ch:analytic-comb]{Chapter 5}, we discuss “Analytic Combinatorics,” applying complex analysis to systematically derive such expansions. 
\end{itemize}
\clearpage

%=====================================================
% Chapter 5: Analytic Combinatorics
%=====================================================
\section{Analytic Combinatorics}

\subsection{Formal Basis}
\begin{itemize}
  \item \textbf{Key Idea}: Translate a \emph{combinatorial definition} into a \emph{generating function} relationship, then apply classical analysis to extract counting or distribution properties.
  \item \textbf{Combinatorial Class}: A set of discrete objects + a size function s.t.\ “\#objects of size $n$” is finite.
  \item \textbf{Unlabelled vs.\ Labelled}:
    \begin{itemize}
      \item \textbf{Unlabelled}: Atoms (of size 1) are indistinguishable.
      \item \textbf{Labelled}: Atoms of size 1 are distinct; permutations matter.
    \end{itemize}
  \item \textbf{OGF} (Ordinary GF) typically for unlabelled classes; \textbf{EGF} (Exponential GF) for labelled classes.
\end{itemize}

%-----------------------------------------------------
\subsection{Symbolic Method for Unlabelled Classes}
\begin{itemize}
  \item \textbf{Basic Operations}:
    \begin{enumerate}
      \item \textbf{Disjoint Union} $(+)$
      \item \textbf{Cartesian Product} $(\times)$
      \item \textbf{Sequence} $\mathrm{SEQ}(\cdot)$
    \end{enumerate}
  \item \textbf{Transfer to OGFs (Theorem~5.1)}:
    \[
      A(z)+B(z)\quad\leftrightarrow\quad A+B;\quad
      A(z)\,B(z)\quad\leftrightarrow\quad A\times B;\quad
      \frac{1}{1 - A(z)}\quad\leftrightarrow\quad \mathrm{SEQ}(A).
    \]
  \item \textbf{Examples}:
    \begin{itemize}
      \item Bitstrings: $B=\mathrm{SEQ}(Z_0+Z_1)\Rightarrow B(z)=\frac{1}{1-2z}.$
      \item No-00 Bitstrings: $G=\varepsilon+(Z_0+Z_1)\times G\;\Longrightarrow\;G(z)=\frac{1+z}{1 - z - z^2}.$
      \item Binary Trees: $T=Z_{\square}+Z_{\bullet}\times T\times T$, etc.
    \end{itemize}
\end{itemize}

%-----------------------------------------------------
\subsection{Symbolic Method for Labelled Classes}
\begin{itemize}
  \item \textbf{Basic Operations (Labelled)}:
    \begin{enumerate}
      \item \textbf{Disjoint Union} $(+)$
      \item \textbf{Star Product} $(\star)$
      \item \textbf{Sequence} $\mathrm{SEQ}(\cdot)$
      \item \textbf{Set} $\mathrm{SET}(\cdot)$
      \item \textbf{Cycle} $\mathrm{CYC}(\cdot)$
    \end{enumerate}
  \item \textbf{Transfer to EGFs (Theorem~5.2)}:
    \[
      A^*(z)+B^*(z)\quad\leftrightarrow\quad A + B;\quad
      A^*(z)\,B^*(z)\quad\leftrightarrow\quad A\star B;
    \]
    \[
      \frac{1}{1 - A^*(z)}\quad\leftrightarrow\quad \mathrm{SEQ}(A);\quad
      e^{A^*(z)}\quad\leftrightarrow\quad \mathrm{SET}(A);\quad
      \ln\frac{1}{1-A^*(z)}\quad\leftrightarrow\quad \mathrm{CYC}(A).
    \]
  \item \textbf{Examples}:
    \begin{itemize}
      \item \emph{Permutations}: $P=\mathrm{SEQ}(Z)\;\Longrightarrow\;P^*(z)=\frac{1}{1-z}.$
      \item \emph{Urns (Sets of distinct atoms)}: $U=\mathrm{SET}(Z)\;\Rightarrow\;U^*(z)=e^z.$
      \item \emph{Cycles}: $C=\mathrm{CYC}(Z)\;\Rightarrow\;C^*(z)=-\ln(1-z).$
    \end{itemize}
\end{itemize}

%-----------------------------------------------------
\subsection{Symbolic Method for Parameters}
\begin{itemize}
  \item \textbf{Bivariate Generating Functions (BGFs)} to track size + cost.
  \item \textbf{Unlabelled} (Ordinary BGF):
    \[
      A(z,u)=\sum_{n,k} a_{n,k}\,z^n\,u^k.
    \]
  \item \textbf{Labelled} (Exponential BGF):
    \[
      A^*(z,u)=\sum_{n,k} a_{n,k}\frac{z^n}{n!}u^k.
    \]
  \item \textbf{Same Transfer Principles (Theorem~5.3 and 5.4)}: Just replace $z$ by $(z,u)$ in the symbolic constructions.
  \item \textbf{Examples}:
    \begin{itemize}
      \item \emph{Bitstrings + number of 1-bits}: $B(z,u)=\frac{1}{1-(1+u)z}.$
      \item \emph{Permutations + \#cycles}: $P^\star(z,u)=\exp\bigl(u\ln\frac{1}{1-z}\bigr),$ etc.
      \item \emph{Binary Trees + \#leaves}, etc.
    \end{itemize}
\end{itemize}

%-----------------------------------------------------
\subsection{Generating Function Coefficient Asymptotics}
\begin{itemize}
  \item \textbf{Goal}: Translate GF expressions (often from the symbolic method) into coefficient approximations.
  \item \textbf{Techniques We’ve Seen}:
    \begin{enumerate}
      \item \emph{Taylor’s Theorem} for exact coefficients of entire functions.
      \item \emph{Rational GF Methods} for poles $\implies$ exponentials or polynomial growth (Theorem~4.1).
      \item \emph{Euler--Maclaurin Summation} to approximate partial sums in expansions.
      \item \emph{Radius of Convergence + Partial Fraction Tools} (Theorem~5.5).
    \end{enumerate}
  \item \textbf{Examples}:
    \begin{itemize}
      \item \emph{Generalized Derangements} ($M$-avoid short cycles). Symbolic approach $\implies$ $f(z)=\exp\bigl(\dots\bigr)$ factor $\implies$ apply Theorem~5.5.
      \item \emph{Catalan Numbers} from $1-\sqrt{1-4z}$ expansions $\implies$ $C_n\sim\frac{4^n}{n^{3/2}\sqrt{\pi}}$.
      \item \emph{2-Regular Graphs}, many other classical enumerations.
    \end{itemize}
  \item \emph{Lesson:} \emph{Singularities} (where the GF “blows up”) determine asymptotic growth of coefficients.
\end{itemize}

%-----------------------------------------------------
% Summary for Chapter 5
%-----------------------------------------------------
\subsection{Summary}
\begin{itemize}
  \item \textbf{Analytic Combinatorics}:
    \begin{enumerate}
      \item \emph{Symbolic Phase}: Define combinatorial classes via algebraic (un)labelled constructions; derive GF relationships.
      \item \emph{Analytic Phase}: Apply real/complex analysis to extract coefficient growth, distribution, etc.
    \end{enumerate}
  \item \textbf{Symbolic Method Transfer Theorems}:
    \begin{itemize}
      \item Unlabelled classes $\Rightarrow$ OGFs using $+$, $\times$, $\mathrm{SEQ}(\cdot)$.
      \item Labelled classes $\Rightarrow$ EGFs using $+$, $\star$, $\mathrm{SEQ}(\cdot)$, $\mathrm{SET}(\cdot)$, $\mathrm{CYC}(\cdot)$.
      \item Bivariate GFs track size + other parameters simultaneously.
    \end{itemize}
  \item \textbf{Analytic Phase} (Transfer to Coefficients):
    \begin{itemize}
      \item Basic expansions, radius-of-convergence arguments, partial fractions, Euler--Maclaurin.
      \item E.g.\ binomial, factorial, Catalan, permutations, derangements, etc.
    \end{itemize}
  \item \emph{Real-World Tip:} Once symbolic GF is obtained, advanced \emph{complex} analysis (singularity analysis) can systematically yield asymptotics for large classes of structures. (See \emph{Analytic Combinatorics} by Flajolet--Sedgewick.)
\end{itemize}
\clearpage
%=====================================================
% Chapter 6: Trees
%=====================================================
\section{Trees}

\subsection{Binary Trees}
\begin{itemize}
  \item \textbf{Definition}: A binary tree is either an external (empty) node or an internal node with exactly two child subtrees (left and right).
  \item \textbf{Catalan Counting}:
    \[
      \text{Number of binary trees with }N\text{ internal nodes} = 
      T_N = \frac{1}{N+1}\binom{2N}{N}.
    \]
  \item \textbf{Useful Properties}:
    \begin{itemize}
      \item \emph{Number of external nodes} = $N+1$ for $N$ internal nodes.
      \item Alternate representations (parentheses, plus-minus path, triangulations).
    \end{itemize}
  \item \textbf{Recursive Formulas}:
    \[
      |t| = |t_L| + |t_R| + 1,\quad
      \mathrm{ipl}(t)=\mathrm{ipl}(t_L)+\mathrm{ipl}(t_R)+|t| - 1,\quad
      h(t)=1+\max\{h(t_L),h(t_R)\}.
    \]
  \item \emph{Actionable Tip:} In many algorithmic contexts, consider $N$ as number of internal nodes or external nodes depending on convenience.
\end{itemize}

%-----------------------------------------------------
\subsection{Forests and Trees}
\begin{itemize}
  \item \textbf{Definition (General)}:
    \begin{itemize}
      \item A \emph{forest} is a \emph{sequence} of disjoint trees.
      \item A \emph{tree} is a root node connected to a forest (i.e.\ the children).
    \end{itemize}
  \item \textbf{Catalan Counting}:
    \[
      \text{Number of forests with }N\text{ nodes} = T_N,\quad
      \text{Number of trees with }N\text{ nodes} = T_{N-1}.
    \]
    i.e.\ $T(z) = 1 + zT(z)$ in OGFs or $T(z)=ze^{T(z)}$ in EGFs for labelled, etc.
  \item \textbf{Correspondences}: 
    \begin{itemize}
      \item \textit{Rotation} / \textit{first-child next-sibling} transforms forests $\leftrightarrow$ binary trees.
      \item Parentheses, gambler's ruin paths, lattice paths $\ldots$
    \end{itemize}
\end{itemize}

%-----------------------------------------------------
\subsection{Combinatorial Equivalences}
\begin{itemize}
  \item \textbf{Rotation Correspondence}: 
    \[
      \text{Forests} \;\longleftrightarrow\;\text{Binary Trees}.
    \]
  \item \textbf{Parenthesis Systems}: Each forest $\leftrightarrow$ Balanced parentheses string.
  \item \textbf{Ruin Sequences} and \textbf{Ballot Problems}: $\pm$ path not crossing zero $\leftrightarrow$ binary tree shape.
  \item \textbf{Lattice Paths} (Dyck paths for Catalan objects) $\leftrightarrow$ binary trees.
  \item \emph{Lesson:} Many classical bijections revolve around Catalan objects (binary trees, partial parentheses, Dyck paths, $\ldots$).
\end{itemize}

%-----------------------------------------------------
\subsection{Properties of Trees}
\begin{itemize}
  \item \textbf{Path Length}:
    \begin{itemize}
      \item Sum of node-levels in a (rooted) tree. 
      \item For binary trees, define \textit{internal} path length and \textit{external} path length.
    \end{itemize}
  \item \textbf{Height}:
    \[
      h(t)=\max\{\text{level of any node}\}.
    \]
  \item \textbf{Recursive definitions} for binary trees $t$:
    \[
      \mathrm{ipl}(t)=\mathrm{ipl}(t_L)+\mathrm{ipl}(t_R)+|t|-1,\quad
      h(t)=1+\max\{h(t_L),h(t_R)\}.
    \]
  \item \textbf{Elementary Bounds}:
    \[
      \mathrm{ipl}(t)\geq \frac{h(t)(h(t)-1)}{2}, \quad
      \mathrm{ipl}(t)\le N\cdot h(t).
    \]
\end{itemize}

%-----------------------------------------------------
\subsection{Examples of Tree Algorithms}
\begin{itemize}
  \item \textbf{Traversal (Pre-/Post-/In-order)}: 
    \begin{itemize}
      \item Recursively visit root and subtrees. 
      \item Stack space usage $\sim$ height.
    \end{itemize}
  \item \textbf{Expression Trees}: 
    \begin{itemize}
      \item Operators in internal nodes, operands in leaves (external). 
      \item Evaluation cost $\sim$ path length or height, or \#registers needed = nontrivial measure.
    \end{itemize}
  \item \emph{Lesson:} The same “shape” of a tree often explains the recursion depth (stack usage) or \#subproblems, etc.
\end{itemize}

%-----------------------------------------------------
\subsection{Binary Search Trees (BSTs)}
\begin{itemize}
  \item \textbf{Definition}: A binary tree in which the key in each node is larger than all keys in its left subtree and smaller than all keys in its right subtree.
  \item \textbf{Insertion Cost}:
    \[
      \mathrm{build\_cost}(\text{tree}) \;=\;\mathrm{ipl}(t)\;\sim\;1.386\,N\log N.
    \]
  \item \textbf{Search Cost}:
    \begin{itemize}
      \item \emph{Successful search}: ~ $2H_N-3$ on average (random permutation). 
      \item \emph{Unsuccessful search}: ~ $2H_{N+1}-2$.
    \end{itemize}
  \item \textbf{Contrast with “Catalan Binary Trees”} (all shapes equally likely). Random permutations $\implies$ BST shapes are more balanced on average.
\end{itemize}

%-----------------------------------------------------
\subsection{Average Path Length in Random Catalan Trees}
\begin{itemize}
  \item \textbf{Two Models}:
    \begin{enumerate}
      \item \emph{Binary Catalan Trees}: each shape of $N$-node tree equally likely.
      \item \emph{BSTs}: shape distribution from random insertions (permutations).
    \end{enumerate}
  \item \textbf{Catalan Path Length}:
    \begin{itemize}
      \item \emph{Binary Trees}: 
        \[
          \mathrm{ipl}_N = 2(N+1)(H_{N+1}-1)-2N \sim 1.356\,N^{3/2}.
        \]
      \item \emph{General Trees}: 
        \[
          \mathrm{pl}_N = (N+1)(2H_N - 2) \sim 1.487\,N^{3/2}.
        \]
    \end{itemize}
\end{itemize}

%-----------------------------------------------------
\subsection{Path Length in Binary Search Trees}
\begin{itemize}
  \item \textbf{Construction Cost} $= \mathrm{ipl}(t)$ for random BST
    \[
      \mathrm{AvgCost} \sim 2(N+1)(H_{N+1}-1)-2N \sim 1.386\,N\log N.
    \]
  \item \textbf{Search Cost}:
    \begin{itemize}
      \item \emph{Successful}: $\sim 2H_N - 3$.
      \item \emph{Unsuccessful}: $\sim 2H_{N+1}-2$.
    \end{itemize}
  \item \emph{Compare to “Random Catalan Trees” shape cost $\sim N^{3/2}$}. 
\end{itemize}

%-----------------------------------------------------
\subsection{Additive Parameters of Random Trees}
\begin{itemize}
  \item \textbf{Additive Parameters} $c(t)$ satisfy $c(t) = e(t) + \sum_{\text{subtrees }s}c(s)$.
  \item \textbf{CGF Approach}:
    \[
      \begin{aligned}
        C_T(z) &=& zE_T(z^?)+\dots & \text{(binary Catalan)}\\
        C_G(z) &=& zE_G(\dots)+\dots & \text{(general Catalan)}\\
        C_B(z) &=& zE_B(\dots)+\dots & \text{(BSTs)}
      \end{aligned}
    \]
  \item \emph{Examples}: \emph{path length}, \emph{\#leaves}, \emph{internal nodes}, etc.
\end{itemize}

%-----------------------------------------------------
\subsection{Height}
\begin{itemize}
  \item \textbf{Non-additive Parameter}. (The maximum of subtree heights, plus 1, not the sum.)
  \item \textbf{Catalan}:
    \begin{itemize}
      \item \textit{Binary Trees}: $\mathrm{height}\sim c\sqrt{N}$, $c\approx 0.95$. 
      \item \textit{General Trees}: $\mathrm{height}\sim c'\sqrt{N}$, $c'\approx 0.765$. 
      \item Harder analysis (Flajolet--Odlyzko, 1982).
    \end{itemize}
  \item \textbf{BSTs}:
    \[
      \mathrm{height}\sim c\log N, \quad c\approx 4.31107,
    \]
    \text{(Devroye, 1986).}
  \item \emph{Non-additivity} often $\implies$ complicated functional eqns + advanced methods (singularity analysis).
\end{itemize}

%-----------------------------------------------------
\subsection{11 Summary of Results}
\begin{itemize}
  \item \textbf{Additive vs Nonadditive}: Path length (easy, linear recursion in GFs) vs Height (harder, “max” of subtrees).
  \item \textbf{Catalan vs BST Models}: Uniform shape vs shape from random permutations $\implies$ different $\mathrm{ipl}$ and $\mathrm{height}$ scales.
  \item \textbf{Overall}, (Catalan) path lengths $\sim N^{3/2}$; (BST) $\sim N\log N$. 
  \item \textbf{Heights}: (Catalan) $\sim c\sqrt{N}$; (BST) $\sim c\log N$.
\end{itemize}

%-----------------------------------------------------
\subsection{Lagrange Inversion}
\begin{itemize}
  \item \textbf{Theorem (Lagrange)}: If $A(z)$ satisfies $z=f(A(z))$, with $f(0)=0$ and $f'(0)\neq 0,$ then
    \[
      [z^n]A(z) = \frac{1}{n}\,[u^{n-1}]\bigl(f(u)\bigr)^n.
    \]
  \item \textbf{Examples}:
    \begin{itemize}
      \item Binary Trees $\to$ $z=T(z)-T(z)^2$ $\implies$ Catalan.
      \item Ternary Trees $\to$ $z=T(z)-T(z)^3$.
      \item $k$-forests, etc.
    \end{itemize}
  \item \emph{Actionable Tip}: Lagrange inversion is a powerful tool for extracting expansions from implicit GF definitions common to tree enumeration problems.
\end{itemize}

%-----------------------------------------------------
\subsection{Rooted Unordered Trees}
\begin{itemize}
  \item \textbf{Definition (Unordered)}: A root node connected to a \textit{multiset} of subtrees.
  \item \textbf{Enumeration} (unlabelled):
    \[
      U(z)=\sum_{n\ge1}U_n\,z^n 
      \quad\text{satisfies}\quad
      \ln(U(z))=\sum_{n\ge1}\frac{U_n}{n}\ln\bigl(1-z^n\bigr),
    \]
    and $U_n\sim c\alpha^n n^{-3/2}$ with $(c,\alpha)\approx (0.4399,2.9558).$
  \item \textbf{Labelled} (Cayley Trees):
    \[
      C(z)=z\,e^{C(z)},\quad
      C_n=n^{\,n-1}.
    \]
\end{itemize}

%-----------------------------------------------------
\subsection{Labelled Trees}
\begin{itemize}
  \item \textbf{Ordered Labelled Trees}: Each node in a fully ordered sequence of subtrees, all node labels distinct $\implies$ $(2N-2)!/(N-1)!$ for $N$-node trees.
  \item \textbf{Cayley (Unordered Labelled) Trees}: $\displaystyle C_n = n^{n-1}$ (famous result by Cayley).
  \item \emph{Lesson:} For labelled, EGF approach enumerates. Often linked to $ze^{T(z)}$ or $z\mapsto ze^z$, etc.
\end{itemize}

%-----------------------------------------------------
\subsection{Other Types of Trees}
\begin{itemize}
  \item \textbf{$t$-ary Trees}: Each node either external or has exactly $t$ subtrees.
    \[
      T(z)=z+(T(z))^t, \quad
      [z^N]T(z)\sim c_t\,\alpha_t^N\,N^{-3/2}.
    \]
  \item \textbf{$t$-restricted Trees}: Each node has up to $t$ subtrees.
    \[
      T(z)=z\bigl(1+T(z)+T(z)^2+\dots+T(z)^t\bigr).
    \]
  \item \textbf{AVL, 2-3, B-trees, Balanced}:
    \begin{itemize}
      \item Suit many searching applications, worst-case $O(\log N)$ path.
      \item Hard to enumerate or get precise random distribution from permutations.
    \end{itemize}
\end{itemize}

%-----------------------------------------------------
% Summary for Chapter 6
%-----------------------------------------------------
\subsection{Summary}
\begin{itemize}
  \item \textbf{Trees} (and forests) are crucial in the analysis of recursive programs and data structures.
  \item \textbf{Major Varieties}: 
    \begin{itemize}
      \item \emph{Unrooted} (free) vs \emph{rooted} vs \emph{ordered} vs \emph{binary}.
      \item \emph{Labelled} vs \emph{Unlabelled}.
      \item \emph{Height} or \emph{path length} differences among distributions (Catalan, BSTs, etc.).
    \end{itemize}
  \item \textbf{Additive Parameters} are approachable with the same recursive + GF methods, while \emph{nonadditive} parameters (height) require more advanced analysis.
  \item \textbf{Analytic Tools}: Lagrange inversion, partial fraction expansions, advanced singularity analysis (particularly for height).
  \item \emph{Forward Reference}: In \hyperref[ch:analytic-comb]{Chapter 5}, we saw how these tree enumerations connect to advanced combinatorial enumerations. We continue exploring further tree-based structures (trees on permutations, tries, etc.) in subsequent chapters.
\end{itemize}
\clearpage

%=====================================================
% Chapter 7: Permutations
%=====================================================
\section{Permutations}

\subsection{Basic Properties of Permutations}
\begin{itemize}
  \item \textbf{Permutation}: An arrangement of $N$ distinct elements, e.g.\ $p_1 p_2 \dots p_N$.
  \item \textbf{Inversions}: Pairs $(i,j)$ with $i < j$ and $p_i > p_j$. The \emph{inversion table} $q_1 q_2 \dots q_N$ satisfies $0 \le q_j < j$ and
    \[
      \mathrm{inv}(p) = \sum_{j=1}^N q_j.
    \]
  \item \textbf{Left-to-right minima (records)}: Indices $i$ such that $p_k > p_i$ for all $k < i$. 
  \item \textbf{Cycles}: Indices $i_1\to i_2\to \dots \to i_t \to i_1$. Derangement: no cycle of length 1.
  \item \textbf{Involution}: $p$ is its own inverse $\Rightarrow$ cycles are 1-cycles or 2-cycles.
  \item \textbf{Runs}: Maximal consecutive (strictly) increasing subsequences. \emph{Rises, falls, peaks, valleys, double-rises, double-falls} measure local monotonic changes.
  \item \textbf{Increasing subsequence}: Subset of indices $i_1 < i_2< \dots < i_k$ with $p_{i_1} < p_{i_2} < \dots < p_{i_k}$.
\end{itemize}

%-----------------------------------------------------
\subsection{Algorithms on Permutations}
\begin{itemize}
  \item \textbf{Sorting Algorithms}: Many methods produce or analyze permutations of keys. 
    \begin{itemize}
      \item \emph{Comparison-based sorting} $\Rightarrow$ random permutations as input model.
    \end{itemize}
  \item \textbf{Rearrangement (Permutation Inversion)}: e.g.\ sorting can compute a permutation $p$, then in-place transform array by applying $p$ (or $p^{-1}$).
  \item \textbf{Randomizing an Array}: Program 7.1: shuffle array to produce random permutation.
  \item \textbf{Priority Queues} / \emph{Heaps}: also interpret permutations as heap-ordered trees.
\end{itemize}

%-----------------------------------------------------
\subsection{Representations of Permutations}
\begin{itemize}
  \item \textbf{Cycle Representation}: $p$ is a set of cycles $(i_1\,i_2\,\dots\,i_t)$.
  \item \textbf{Foata's Correspondence}: Minimizing each cycle, cycle leaders in decreasing order $\leftrightarrow$ left-to-right minima.
  \item \textbf{Inversion Tables}: $q_1 \dots q_N$ with $0 \le q_i < i$. 1-to-1 with permutations.
  \item \textbf{Lattice Representation}: Mark $(p_i,i)$ in an $N\times N$ grid. Inversions correspond to cells with row>row-of$(i)$, column<column-of$(i)$.
  \item \textbf{BST Representation}: Insert $p_1,\dots,p_N$ successively. (Many permutations map to the same BST shape.)
  \item \textbf{HOT (Heap-Ordered Tree) Representation}: Insert $p_i$ in a min-heap shape. Each permutation has a unique HOT.
\end{itemize}

%-----------------------------------------------------
\subsection{Enumeration Problems}
\begin{itemize}
  \item \textbf{Cycle-length constraints}: 
    \begin{itemize}
      \item $\mathrm{EGF}=$ product expansions from the symbolic method: 
        \[
          \exp\!\Bigl(\frac{z^k}{k}\Bigr) \quad\text{(cycles of length }k\text{)}.
        \]
      \item \emph{Derangements}: no fixed points (1-cycles) $\Rightarrow$ EGF $= \exp(z+z^2/2 +\dots)$ etc.
    \end{itemize}
  \item \textbf{Involutions}: cycles of only lengths 1 or 2. $\mathrm{EGF}= e^{z + z^2/2}$. 
  \item \textbf{Maximal/Minimal cycle length} enumerations: partial sums give $\mathrm{EGF}$ via the cycle restrictions. 
\end{itemize}

%-----------------------------------------------------
\subsection{Analyzing Properties with CGFs}
\begin{itemize}
  \item \textbf{Additive vs Nonadditive}: As with trees, sums are easier (like \#inversions, \#left-to-right minima); maxima (\# passes) are harder.
  \item \textbf{General Tools}: define $\displaystyle B(z)= \sum_{p} (\text{cost}(p))\,\frac{z^{|p|}}{|p|!}$, exploit combinatorial construction to get functional/differential eqn.
\end{itemize}

%-----------------------------------------------------
\subsection{Inversions and Insertion Sort}
\begin{itemize}
  \item \textbf{Insertion Sort}: cost $\sim$ \#inversions in permutation.
  \item \textbf{Distribution of Inversions}: 
    \begin{itemize}
      \item PGF factorization: each $q_i$ uniform in $[0\dots i-1]$ $\implies$ 
      \[
        \sum_{k} \Pr\{\mathrm{inv}=k\}u^k 
        = \prod_{i=1}^N \bigl(1+u+u^2+\dots+u^{i-1}\bigr)/i.
      \]
      \item Mean $\mathrm{inv}\approx \frac{N(N-1)}{4}$, stdev $\approx \sqrt{\frac{N(N-1)(2N+5)}{72}}$.
    \end{itemize}
  \item \textbf{EGF approach}:
    \[
      B(z)= \sum \mathrm{inv}(p)\,\frac{z^N}{N!}
      = \frac{z^2}{2(1-z)^3}.
    \]
  \item \textbf{Conclusion}: insertion sort $\sim \frac{N^2}{4}$ compares/moves on average.
\end{itemize}

%-----------------------------------------------------
\subsection{Left-to-Right Minima and Selection Sort}
\begin{itemize}
  \item \textbf{Left-to-right minima} $\leftrightarrow$ \textbf{cycles} by Foata's representation.
  \item \textbf{Stirling \#'s of first kind}: 
    \[
      \mathrm{EGF} \quad B(z,u)=(1-z)^{-u}.
    \]
    The number of permutations with $k$ left-to-right minima is $\left\lvert {n \brack k} \right\rvert$, the unsigned Stirling cycle numbers.
  \item \textbf{Selection Sort}: cost mostly from minimum searching, $\sim N^2/2$ compares + $N - 1$ swaps. 
    \begin{itemize}
      \item The $\mathrm{lrm}(p)$ updates is $\sum \mathrm{lrm}(p)= (N+1)(H_{N+1}-1)$.
      \item $\sim N \ln N$ updates (dominates).
    \end{itemize}
\end{itemize}

%-----------------------------------------------------
\subsection{Cycles and In Situ Permutation}
\begin{itemize}
  \item \textbf{Permutation Rewriting In Place}: Program 7.5 uses cycles, each cycle is rotated once $\implies$ \#moves $= N - \#1\text{-cycles}$, \#outer loops $= \#\text{cycles}$.
  \item \textbf{Average \#cycles = } $H_N$. \quad Average \#1-cycles $\approx 1$.
  \item \textbf{Cycle length distribution}:
    \begin{itemize}
      \item Probability that $\mathrm{maxCycleLength}\le k$ from enumerations in \S 7.4.
    \end{itemize}
\end{itemize}

%-----------------------------------------------------
\subsection{Extremal Parameters}
\begin{itemize}
  \item \textbf{Bubble Sort}:
    \begin{itemize}
      \item \#exchanges $=$ \#inversions,
      \item \#passes $=$ $\max\{\text{inversion table entry}\} \sim \frac{\ln N}{\ln\ln N}$ (Ramanujan Q function). 
    \end{itemize}
  \item \textbf{Longest Cycle} in a random perm $\sim 0.6243\,N$. 
  \item \textbf{Longest run} or \textbf{longest increasing subsequence}: More advanced results, typically via advanced combinatorial/probabilistic or analytic arguments (e.g.\ $\sim 2\sqrt{N}$ for LIS).
\end{itemize}

%=====================================================
% Summary
%=====================================================
\subsection{Summary}
\begin{itemize}
  \item \textbf{Permutations} are fundamental combinatorial objects with many properties relevant to sorting and rearrangement algorithms.
  \item \textbf{Properties}: inversions, left-to-right minima, cycles, runs, etc. Each has direct interpretation in an algorithm (insertion sort, selection sort, in situ permutation, bubble sort, etc.).
  \item \textbf{Analysis Tools}: 
    \begin{itemize}
      \item \emph{CGFs} (cumulative generating functions) often yield average values with linear recurrences.
      \item \emph{PGFs} for distributions (like inversions) factor nicely.
      \item \emph{EGFs} with symbolic method for cycle constraints (e.g.\ derangements, involutions).
    \end{itemize}
  \item \textbf{Extremal parameters} (longest cycle, \# passes in bubble sort) are typically harder to analyze; advanced methods (singularity analysis) often required.
  \item \emph{Applications}:
    \begin{itemize}
      \item Insertion sort $\sim \frac{N^2}{4}$, selection sort $\sim \frac{N^2}{2}$, bubble sort $\sim \frac{N^2}{2}$ compares, plus subtle pass counts.
      \item Many open problems (shellsort) demonstrate the complexity behind “simple” permutation algorithms.
    \end{itemize}
\end{itemize}
\clearpage

%=====================================================
% Chapter 8: Strings and Tries
%=====================================================
\section{Strings and Tries}

\subsection{String Searching}
\begin{itemize}
  \item \textbf{Problem setup:} Given a pattern of length $P$ and a text of length $N$, find all (or the first) occurrences of the pattern in the text.
  \item \textbf{Basic (Naive) Algorithm:} For each position in the text (0 to $N-P$), check the next $P$ characters for a match.
  \item \textbf{Cost analysis:} 
    \begin{itemize}
      \item Occurrences of \emph{pattern prefixes} at each text position.
      \item Expected cost $\approx 2N$ bit inspections for random text/pattern, but can degrade with highly repetitive patterns (e.g.\ many zeros).
    \end{itemize}
\end{itemize}

%-----------------------------------------------------
\subsection{Combinatorial Properties of Bitstrings}
\begin{itemize}
  \item \textbf{Bitstring model:} Each of the $2^N$ bitstrings of length $N$ equally likely.
  \item \textbf{Classical Bernoulli trials:} String properties $\leftrightarrow$ local or global events in coin flips.
  \item \textbf{Pattern occurrences:} 
    \begin{itemize}
      \item Position of first occurrence $\leftrightarrow$ enumerating bitstrings that avoid the pattern.
      \item Counting uses GF translations of set decompositions.
    \end{itemize}
  \item \textbf{Runs of 0s:} 
    \begin{itemize}
      \item OGF for strings avoiding $P$ consecutive zeros:
        \[
          B_P(z) = \frac{(1-z^P)}{1-z - z(1-z^P)} \quad\text{(or a variant).}
        \]
      \item Average position of first run of $P$ zeros is $\approx 2^P$ in random bitstrings.
      \item Probability that $P$ zeros do not occur is exponentially small in $N$ for fixed $P$.
    \end{itemize}
  \item \textbf{Longest run:} Typically $\approx \log_2(N)$ in a random bitstring, with oscillatory/fluctuation phenomena.
\end{itemize}

%-----------------------------------------------------
\subsection{Regular Expressions}
\begin{itemize}
  \item \textbf{REs:} Formal language specification using union $(+)$, concatenation, and Kleene star $(^*)$.
  \item \textbf{Rational OGFs:} Sets of bitstrings described by unambiguous REs have rational generating functions for enumeration.
  \item \textbf{Examples:} 
    \begin{itemize}
      \item $\{0,1\}^*$ all bitstrings, GF $= 1/(1-2z)$.
      \item No runs of length $\ge P$: see above.
      \item Gambler's ruin sequences as RE $\Rightarrow$ GFs are rational.
    \end{itemize}
\end{itemize}

%-----------------------------------------------------
\subsection{Finite-State Automata and the KMP Algorithm}
\begin{itemize}
  \item \textbf{FSAs:} Recognize languages from REs; can be used for pattern matching.
  \item \textbf{Knuth-Morris-Pratt (KMP):}
    \begin{itemize}
      \item Builds an FSA for a pattern of length $P$ in $O(P)$ time.
      \item Then checks text of length $N$ with only $O(N)$ comparisons.
    \end{itemize}
  \item \textbf{Practical Impact:} Avoids re-checking overlaps for repetitive patterns.
\end{itemize}

%-----------------------------------------------------
\subsection{Context-Free Grammars}
\begin{itemize}
  \item \textbf{CFGs:} More powerful than REs; can describe languages like $\{0^n\,1^n\}$.
  \item \textbf{OGFs:} For unambiguous CFGs, the enumerating GFs are \emph{algebraic} (solutions of polynomial equations).
  \item \textbf{Examples:} 
    \begin{itemize}
      \item Balanced parentheses, Dyck paths, ballot problem, etc.
      \item $n$ zeroes, $n$ ones $\to$ $O(z)$ satisfies polynomial equations. (2-ordered permutations, etc.)
    \end{itemize}
\end{itemize}

%-----------------------------------------------------
\subsection{Tries}
\begin{itemize}
  \item \textbf{Definition:} Combinatorial structure for a set of bitstrings:
    \begin{itemize}
      \item Each external node can represent a bitstring if prefix-free or if marking final bits.
      \item Internal nodes branch by 0 vs.\ 1.
    \end{itemize}
  \item \textbf{Prefix code property:} No stored string is a prefix of another. 
  \item \textbf{Representation:} The path from root to external node identifies each bitstring in the set.
\end{itemize}

%-----------------------------------------------------
\subsection{Trie Algorithms}
\begin{itemize}
  \item \textbf{Trie-based search:} Symbol table with bit-inspections (digital searching).
  \item \textbf{Patricia tries:} Compress one-way branching for memory efficiency.
  \item \textbf{Radix-exchange sort:} Recursively partition by leading bit, cost $\sim N\lg N$ for random bits.
  \item \textbf{Prefix-free codes (Huffman):} External nodes labeled, decoding by traversing trie from root.
  \item \textbf{Suffix tries:} Preprocess text for all pattern searches (cost $\sim N\lg N$).
  \item \textbf{Multiple patterns:} Build trie or FSA (Aho-Corasick) for sets of patterns.
  \item \textbf{Leader election protocol:} Bit-based elimination among $N$ participants also yields trie structure.
\end{itemize}

%-----------------------------------------------------
\subsection{Combinatorial Properties of Tries}
\begin{itemize}
  \item \textbf{Random tries:} Built from $N$ random infinite bitstrings.
  \item \textbf{Path length:} $\mathrm{E}(\text{external path length}) \sim N\log_2 N$ with periodic fluctuations.
  \item \textbf{Number of internal nodes:} $\sim \frac{N}{\ln 2}$ with small oscillations.
  \item \textbf{Algorithms:} 
    \begin{itemize}
      \item Searching in tries $\to O(\log N)$ bit inspections for random data.
      \item Radix-exchange $\to N\log N$ cost.
      \item Leader election $\to \log N$ rounds on average; success probability $\approx 0.72135$.
    \end{itemize}
  \item \textbf{Advanced methods:} Mellin transform, complex analysis, and other analytic combinatorics tools are typically needed to fully capture oscillatory terms.
\end{itemize}

%-----------------------------------------------------
\subsection{Larger Alphabets}
\begin{itemize}
  \item \textbf{General M-ary strings:} Direct analogs of results for binary strings, but constants differ by factors of $\ln M$.
  \item \textbf{Multiway tries:} Factor-of-$\log M$ savings in search cost vs.\ binary tries, at expense of bigger node array (size $M$).
  \item \textbf{Applications:} Ternary tries (TSTs), LZW data compression, etc.
\end{itemize}

%=====================================================
% Summary
%=====================================================
\subsection{Summary}
\begin{itemize}
  \item \textbf{Bitstrings} are fundamental structures relating to Bernoulli processes, pattern occurrences, and local ``runs.''
  \item \textbf{Analysis of basic string-searching algorithms} (naive, KMP, Aho-Corasick, suffix tries) shows that advanced preprocessing can reduce or eliminate redundant comparisons.
  \item \textbf{Formal languages perspective:} 
    \begin{itemize}
      \item \emph{Regular expressions (REs) $\to$ rational generating functions}
      \item \emph{Context-free grammars (CFGs) $\to$ algebraic generating functions}
    \end{itemize}
  \item \textbf{Tries} provide data structures for sets of bitstrings:
    \begin{itemize}
      \item \emph{Digital search tries, Patricia tries, Radix-exchange, Huffman codes, Suffix tries, Multiple-pattern matching, Leader election.}
      \item \emph{Average path length} $\sim N \log_2 N$; number of internal nodes $\sim \frac{N}{\ln 2}$; both exhibit small oscillations.
    \end{itemize}
  \item \emph{Advanced Topics}: exact expansions often require complex analytic methods (Mellin transforms, etc.).
\end{itemize}
\clearpage

%=====================================================
% Chapter 9: Words and Mappings
%=====================================================
\section{Words and Mappings}

\subsection{Hashing with Separate Chaining}
\begin{itemize}
  \item \textbf{Basic idea}: Distribute $N$ keys into $M$ lists (``urns''), where a \emph{hash function} maps each key to an integer between $1$ and $M$. If collisions occur (two or more keys hashing to same value), maintain them in a small secondary structure (linked list).
  \item \textbf{Performance considerations}:
    \begin{itemize}
      \item Typically need $M \approx \alpha N$ ($\alpha$ a constant factor) to keep average chain length $\approx \alpha$.
      \item Simple to implement; widely used in practical applications.
      \item Analyzed via balls-and-urns model.
    \end{itemize}
  \item \textbf{Summary of cost}:
    \begin{itemize}
      \item Unsuccessful search: $\frac{N}{M}$ probes on average.
      \item Successful search: $\frac{N+1}{2M}$ probes on average.
    \end{itemize}
\end{itemize}

%-----------------------------------------------------
\subsection{The Balls-and-Urns Model and Properties of Words}
\begin{itemize}
  \item \textbf{Definition (words)}: An $M$-word of length $N$ is a function $f: \{1,\dots,N\}\to\{1,\dots,M\}$. Equivalently, a sequence of length $N$ over an alphabet of size $M$.
  \item \textbf{Balls-and-urns model}:
    \begin{itemize}
      \item Distributing $N$ labelled balls into $M$ labelled urns.
      \item Widely studied classical combinatorial model; many applications in hashing, occupancy, random distributions.
    \end{itemize}
  \item \textbf{Occupancy distribution}:
    \begin{itemize}
      \item Probability that $k$ balls land in a particular urn is given by the binomial distribution $\binom{N}{k} \left(\frac{1}{M}\right)^{k}\!\left(\frac{M-1}{M}\right)^{N-k}$.
      \item Key structural parameters: number of empty urns, max/min occupancy, typical occupancy $N/M$, etc.
    \end{itemize}
\end{itemize}

%-----------------------------------------------------
\subsection{Birthday Paradox and Coupon Collector Problem}
\begin{itemize}
  \item \textbf{Birthday problem}:
    \begin{itemize}
      \item How many balls need to be thrown before some urn (hash value) contains $2$ balls for the first time?
      \item Probability that no collision occurs in $N$ throws is $\frac{(M)_N}{M^N}$ (falling factorial); typically $\approx e^{-N^2/(2M)}$ for large $N<M$.
      \item Number of people needed so the chance of a shared birthday $>1/2$ is about $1.2\sqrt{M}$ when $M=365$ (roughly $23$ people).
    \end{itemize}
  \item \textbf{Coupon collector problem}:
    \begin{itemize}
      \item How many balls must be thrown so that all urns are nonempty at least once?
      \item Average $\approx M\ln(M)$; more precisely, $M(H_M) = M\ln M + \gamma M + O(1)$.
    \end{itemize}
\end{itemize}

%-----------------------------------------------------
\subsection{Occupancy Restrictions and Extremal Parameters}
\begin{itemize}
  \item \textbf{Maximal occupancy constraints}:
    \begin{itemize}
      \item EGF for words with at most $k$ copies of each letter: $(1 + z + z^2/2! + \cdots + z^k/k!)^M$.
      \item Generalization of birthday problem to having at most $k$ balls in each urn.
    \end{itemize}
  \item \textbf{Minimal occupancy constraints}:
    \begin{itemize}
      \item EGF for words with at least $k$ copies of each letter: $\bigl(e^z - 1 - z - \cdots - z^k/k!\bigr)^M$.
      \item Generalization of coupon collector to having at least $k$ balls in each urn.
    \end{itemize}
  \item \textbf{Max/min occupancy extremes}:
    \begin{itemize}
      \item Largest number of balls in any urn (average grows like $\frac{\ln N}{\ln \ln N}$).
      \item Number of urns with occupancy $=0$, $=1$, etc., described by binomial/Poisson approximations.
    \end{itemize}
\end{itemize}

%-----------------------------------------------------
\subsection{Occupancy Distributions}
\begin{itemize}
  \item \textbf{Binomial distribution}:
    \[
      \Pr\{\text{urn has }k\text{ balls}\}
      \;=\;
      \binom{N}{k} \left(\frac{1}{M}\right)^k \left(\frac{M-1}{M}\right)^{N-k}.
    \]
  \item \textbf{Poisson approximation} when $N$ grows, $M$ also grows with $N$ (load factor $\alpha = N/M$ fixed or $M=\lambda N$):
    \[
      \Pr\{k\} \approx \frac{\alpha^k e^{-\alpha}}{k!}.
    \]
  \item \textbf{Normal approximation} if $M$ small relative to $N$, or focusing on large occupancy counts near $N/M$.
  \item \textbf{Example (hashing with separate chaining)}:
    \begin{itemize}
      \item Number of probes for unsuccessful search $\approx N/M$,
      \item For successful search $\approx (N+1)/(2M)$.
    \end{itemize}
\end{itemize}

%-----------------------------------------------------
\subsection{Open Addressing Hashing}
\begin{itemize}
  \item \textbf{Idea}: Store all keys in a single array; collisions resolved by finding next open slot. Reduces pointer overhead but can cause ``clustering.''
  \item \textbf{Linear probing}:
    \begin{itemize}
      \item Insert key at hashed position; if occupied, move left until finding empty slot.
      \item Average search cost (for load factor $\alpha$):
        \[
          \approx \frac{1}{2}\left(\frac{1}{1-\alpha}\right)+\frac{1}{2}
          \quad\text{successful,}
          \quad
          \approx \frac{1}{(1-\alpha)^2}
          \quad\text{unsuccessful.}
        \]
      \item Degrades badly as $\alpha\to1$ (long clusters).
    \end{itemize}
  \item \textbf{Uniform hashing}:
    \begin{itemize}
      \item Assumes each permutation of table positions equally likely.
      \item Successful search $\approx \frac{1}{1-\alpha}$; unsuccessful $\approx \frac{1}{(1-\alpha)^2}$.
      \item Achieves smaller cost than linear probing, especially when table nearly full.
    \end{itemize}
\end{itemize}

%-----------------------------------------------------
\subsection{Mappings}
\begin{itemize}
  \item \textbf{Definition}: $N$-mapping is a function $f: \{1,\dots,N\}\to\{1,\dots,N\}$. There are $N^N$ such mappings.
  \item \textbf{Graph structure}:
    \begin{itemize}
      \item Each node has outdegree 1, possibly many in-degrees.
      \item Decomposes into cycles of nodes, each with trees attached (a \emph{tail+cycle} or ``rho'' structure).
    \end{itemize}
  \item \textbf{Statistics of interest}:
    \begin{itemize}
      \item Number of cycles, size of each cycle, number of nodes on cycles, length of tails, etc.
      \item Symbolic method: consider ``sets of cycles of trees.'' 
      \[
        \text{EGF} = \exp\!\Bigl(\log\bigl(\frac{1}{1-C(z)}\bigr)\Bigr) = \dots
      \]
    \end{itemize}
  \item \textbf{Key results} (for random $N$-mappings):
    \begin{itemize}
      \item Number of components: $\sim 0.624\dots N$ on average.
      \item Number of nodes on cycles: $\sim 0.757\dots N$ on average.
      \item Rho length of a random point: $\sim c$ (the Ramanujan $Q$-function $\approx 2.606\dots$).
      \item Largest cycle $\sim c' N^{1/2}$ or $\sim \frac{\ln N}{\ln\ln N}$ depending on regime (details from advanced references).
    \end{itemize}
\end{itemize}

%-----------------------------------------------------
\subsection{Integer Factorization and Mappings}
\begin{itemize}
  \item \textbf{Pollard's rho method}:
    \begin{itemize}
      \item Exploits random mapping structure to find nontrivial factors of $N$ (assume $N = pq$).
      \item Uses iteration of $f(x) = x^2 + c\mod N$ (like a random mapping).
      \item Floyd's cycle-finding algorithm (constant space, $O(\sqrt{p})$ time) locates collisions mod $p$ before collisions mod $q$ if $p < q$.
    \end{itemize}
  \item \textbf{Heuristic analysis}: If $f(x)$ behaves like a random $N$-mapping, then iteration length (the rho length) is $O(\sqrt{p})$, thus $O(N^{1/4})$ steps for $N = pq$. 
  \item \textbf{Practical performance}: Efficient for 10--30 digit $N$; pollard's method discovered factor of Fermat number $F_8$ in 1980.
\end{itemize}

%=====================================================
% Chapter 9 Summary
%=====================================================
\subsection{Summary}
\begin{itemize}
  \item \textbf{Words (M-ary strings)} relate to Bernoulli trials and balls-and-urns occupancies.
  \item \textbf{Hashing with separate chaining} analyzed via occupancy distribution:
    \begin{itemize}
      \item Unsuccessful search: $\approx N/M$ comparisons,
      \item Successful search: $\approx (N+1)/(2M)$ comparisons.
    \end{itemize}
  \item \textbf{Classical problems}: 
    \begin{itemize}
      \item \emph{Birthday paradox}: first collision after about $\sqrt{M}$ tries,
      \item \emph{Coupon collector}: fill $M$ urns in about $M\ln M$ tries.
    \end{itemize}
  \item \textbf{Open addressing hashing} (linear probing, uniform hashing) studied in detail; cost depends on load factor $\alpha = N/M$.
  \item \textbf{Mappings}: generalize permutations and trees, each node has outdegree $=1$:
    \begin{itemize}
      \item Decompose into cycles and trees,
      \item Symbolic method yields enumerations and average properties,
      \item $N^N$ total; each mapping has a number of \emph{components}, \emph{cycle lengths}, \emph{tail lengths}, etc.
    \end{itemize}
  \item \textbf{Pollard's rho method} for integer factorization exemplifies the use of random mappings and cycle-finding (Floyd's algorithm). Performance $\approx O(N^{1/4})$ on average under standard heuristic assumptions.
\end{itemize}
% -------------------------------------------------------------
% End of Cheat Sheet
% -------------------------------------------------------------

\end{document}

